\section{Composition with the deployment-safety theorem}
\label{sec:composition}

Theorem~\ref{thm:c7-corollary} (bounded co-evolution corollary)
and Theorem~\ref{thm:concgap-scoped} (scoped Concentration-Gap)
each enter the deployment-safety theorem
\citep[Theorem~1]{lasser2026paper10} at specific points.  This
section locates those points and lists the conditions each
theorem adds to or removes from the deployment-claim hypothesis
set.

\subsection{Conditions removed and added}

\paragraph{Removed.}  The deployment-claim hypothesis set, after
the present theorems, no longer contains:
\begin{itemize}
\item The free-floating bounded co-evolution assumption that
  $\bar{M}$ is bounded by some $|P|$-independent constant.
  Replaced by Theorem~\ref{thm:c7-corollary}.  In
  \citep[Deployment Safety]{lasser2026paper10}, the property is
  now stated as condition~(C7) and invokes the corollary form
  directly.
\item The HHI surrogate-adequacy assumption: the
  converse-direction sufficiency claim $\HHI < H^*
  \Rightarrow \mathcal{R}_{\mathrm{press}}^c$. Replaced by
  Theorem~\ref{thm:concgap-scoped} under the conjunction (REP) +
  (DISP) + (COAL) and embedding compatibility
  (Assumption~\ref{ass:embedding-companion}), now collected as
  \citep[Assumption~1]{lasser2026paper10} (Concentration-Gap
  structural conditions).
\end{itemize}

\paragraph{Added.} The deployment-claim hypothesis set acquires
five explicit hypotheses, each operationally auditable (with
embedding compatibility certified as a precondition for the
embedded representations used by (REP) and (DISP)):
\begin{itemize}
\item \textbf{(C7.RATE).} Upper exposure-rate cap from
  Assumption~\ref{ass:c7-rate}. Added as a sub-clause of
  \citep[Condition~(C7)]{lasser2026paper10}.
\item \textbf{(REP).} Representativeness from
  Definition~\ref{def:rep-companion}.
\item \textbf{(DISP).} Bounded dispersion from
  Definition~\ref{def:disp-companion}.
\item \textbf{(COAL).} Coalition closure from
  Definition~\ref{def:coal-companion}, with explicit residual
  $\eta_{\mathrm{latent}}$ for undetected coordination and
  post-partition counterparty cardinality $N$ deployment-class
  bounded (clause (ii); structurally required for the
  HHI-to-$\chi^2$ translation to be $|P|$-independent).
\item \textbf{Embedding compatibility} from
  Assumption~\ref{ass:embedding-companion}: an affine isometry
  (or bi-Lipschitz, with constants tracked through the bound)
  $\phi: \mathcal{V}_\mathrm{utility} \to \mathcal{V}$ on the
  operationally active subspace. This is the structural
  prerequisite for transferring the algebraic distortion bound to
  $\|\Pproxy - \Ttruth\|$ and thence to Goodhart slack.
\end{itemize}

(REP), (DISP), (COAL), and embedding compatibility are collected
as the structural conditions for the Concentration-Gap selection
bound; (C7.RATE) is the exposure-rate condition for the bounded
co-evolution corollary.

\subsection{Where each theorem enters Deployment Safety's proof}

\paragraph{Theorem~\ref{thm:c7-corollary} enters at Lemma~1's co-evolution substitution.}
\citep[Lemma~1, Step~4]{lasser2026paper10} substitutes the
co-evolution correction term $(\epsnonresCoev)_+ \leq
K_{\mathrm{coev}} \cdot \bar{M}$ into the four-channel
decomposition; \citep[Lemma~1, Step~6]{lasser2026paper10}
packages the result via the generic-$X$ intensivity argument.
Both rely on $\bar{M}$ being a $|P|$-independent constant. The
corollary supplies $\bar{M}^{\mathrm{cum}} \leq C \cdot \Bclip
\cdot \lammax \cdot \taumeta$ with all factors deployment-class
constants intensive in $|P|$, closing the substitution at Step~4
without a free-floating constant.

\paragraph{Theorem~\ref{thm:concgap-scoped} enters at Layer~1 binding via $I_5$.}
\citep[Theorem~1, Layer~1 proof, Step~4]{lasser2026paper10} requires
the converse-direction sufficiency $\HHI < H^* \Rightarrow$
deployment outside $\mathcal{R}_{\mathrm{press}}$. The scoped
selection theorem supplies the Goodhart-slack upper bound
$|g(\Ttruth) - g(\Pproxy)| \leq \Lip(g) \cdot (\rho_\mathrm{rep} +
\sigma \sqrt{\chi^2(w \|\mu)} + \eta_\mathrm{latent})$ under (REP)
+ (DISP) + (COAL) and embedding compatibility
(Assumption~\ref{ass:embedding-companion}); for uniform $\mu$ on
$N$ counterparties,
$\chi^2 = N \HHI - 1$, so $\HHI < H^*$ yields a deployment-class
constant ceiling on the slack.

\paragraph{Both theorems enter Deployment Safety's sensitivity analysis.}
\citep[\S 10.2]{lasser2026paper10}'s failure-mode taxonomy
decomposes into specific structural-condition failures:
bounded-co-evolution failure factors through (C5.HOEFF) /
(C7.RATE) / channel-projection-structure failures (each with
audit-detectable signal), and Concentration-Gap structural-condition
failure factors through (REP) / (DISP) / (COAL) /
embedding-compatibility failures (the last via the bi-Lipschitz
$\phi$ certification of \S\ref{sec:audits}).

\subsection{Deployment-claim shape}

The structural shape of the deployment-safety theorem is
preserved:
\begin{itemize}
\item Eleven operational invariants $I_1, \ldots, I_{11}$.
\item Three-layer claim (static safe region /
  detection-and-correction / acknowledged residuals).
\item Five named residuals (R1)--(R5).
\item Cooperative-anchoring property and canonical tripartite
  substrate identification.
\item Operational conditions (C2)--(C6) and (C8)--(C12) of
  \citep[\S 5]{lasser2026paper10} retain their content;
  (C1) and (C7) carry refined structural dependencies: (C1) now
  explicitly invokes (REP), (DISP), (COAL), and embedding
  compatibility via \citep[Assumption~1]{lasser2026paper10}, and
  (C7) carries the (C7.RATE) sub-clause established here.
\end{itemize}

Theorem~\ref{thm:c7-corollary} and
Theorem~\ref{thm:concgap-scoped} close structural gaps in the
deployment-claim hypothesis set without altering the deployment
claim's shape.
